\subsection*{B.2 \quad Meaning of the dual variables}

We show each price as a gain of the margin when the right-hand side grows by one unit; this is $-\pi$ of B.1 (for every row;\texttt{src/print\_b.jl}). Let $\alpha_b=-\pi_1$ be the price of land, $\beta_t=-\pi_5$ of water in the reservoir, $\gamma_t=-\pi_{8a}$ of a family hour.

\begin{figure}[!ht]
\centering
\begin{tikzpicture}
\begin{groupplot}[group style={group size=3 by 1, horizontal sep=1.4cm}, width=5.4cm, height=3.7cm, ymin=0, enlarge x limits=0.2, nodes near coords, nodes near coords style={font=\tiny}]
\nextgroupplot[ybar, bar width=10pt, symbolic x coords={B1,B2,B3,B4}, xtick=data, title={land $\alpha_b$ (EUR/ha)}, ymax=980, every node near coord/.append style={/pgf/number format/.cd,fixed,precision=0}]
\addplot[fill=cmaize] coordinates {(B1,793.01)(B2,641.01)(B3,563.01)(B4,809.01)};
\nextgroupplot[ybar, bar width=10pt, symbolic x coords={P1,P2,P3,P4}, xtick=data, title={water $\beta_t$ (EUR/m$^3$)}, ymax=0.9, every node near coord/.append style={/pgf/number format/.cd,fixed,precision=2}]
\addplot[fill=cbarley] coordinates {(P1,0)(P2,0.69)(P3,0.69)(P4,0.69)};
\nextgroupplot[ybar, bar width=10pt, symbolic x coords={P1,P2,P3,P4}, xtick=data, title={family hour $\gamma_t$ (EUR/h)}, ymax=46, every node near coord/.append style={/pgf/number format/.cd,fixed,precision=1}]
\addplot[fill=cwheat] coordinates {(P1,0)(P2,0)(P3,0)(P4,36.40)};
\end{groupplot}
\end{tikzpicture}
\caption{Dual variables of the land, reservoir and labour constraints.}
\end{figure}


\begin{itemize}
\item \textbf{Land} $\alpha_b$ (EUR/ha): one more hectare is worth 793, 641, 563, 809 EUR on B1--B4.\item \textbf{Water} $\beta_t$ (EUR/m$^3$): 0 in P1, 0.69 in P2--P4. This is $0.60$ (price of the permit) $+\,0.09$ (pumping cost): an extra m$^3$ saves one pumped m$^3$.
\item \textbf{Labour} $\gamma_t$ (EUR/h): 0 in P1--P3, 36.40 in autumn. This is the wage $24$ plus $12.40$ for the hiring cap, which is reached.
\end{itemize}

\textbf{Why is water worth 0 in P1?} The reservoir is full and 1\,500 m$^3$ overflow, so one more m$^3$ in spring is lost. From P2 on, stored water replaces pumped water, which is limited by the permit, so it has a value.
